\documentclass[11pt]{article}
\usepackage[margin=1in]{geometry}
\usepackage{amsmath,amssymb}
\usepackage{booktabs}
\usepackage{graphicx}
\usepackage{float}
\usepackage{hyperref}
\usepackage{xcolor}
\hypersetup{colorlinks=true,linkcolor=blue!50!black,urlcolor=blue!50!black,citecolor=blue!50!black}

\title{\textbf{PE-RankFormer: Models and Results}\\[0.4em]
\large A technical report on the current best system for prime-editing efficiency prediction}
\author{}
\date{22 August 2026}

\begin{document}
\maketitle

\begin{abstract}
\noindent
We report the current best models for predicting prime-editing efficiency, benchmarked
against OptiPrime on its own training data, cross-validation splits, and complete
held-out test set. The final system reaches Spearman $\rho = \mathbf{0.9079}$ on the
20{,}509-row official held-out set against OptiPrime's $0.8690$, a difference of
$\mathbf{+0.0389}$ with a protospacer-clustered paired bootstrap 95\% CI of
$[+0.0286, +0.0498]$, ahead in \emph{all} 5{,}000 resamples ($p < 0.0002$). The margin
is $+0.0804$ on the Kim partition and $+0.0220$ on Liu.

The system is an equal-weight rank average of five out-of-fold members. Its strongest
single member --- an \emph{ordinal-SSM} model combining a cumulative-threshold
objective with a bidirectional state-space sequence mixer --- reaches $0.9082$
out-of-fold on development folds, above the entire previous three-member ensemble.

We also report, in equal detail, what did \emph{not} work. Approximately fifty
subsequent experiments spanning supervision geometry, architecture, capacity, output
parameterisation and training scheme produced no improvement over this system: each
returned what an unmodified re-run of the same model returns. We give the measurements
that establish this, a replicate-based estimate showing that roughly $+0.12$ Spearman
of \emph{genuine} headroom nonetheless remains on the Kim partition, and a diagnosis of
where that headroom sits.
\end{abstract}

\tableofcontents

% ==========================================================================
\section{Summary of results}
% ==========================================================================

\begin{table}[H]
\centering
\caption{Final held-out performance. Official OptiPrime held-out set, $n=20{,}509$
rows over 750 protospacer clusters. No model in this table was trained on these rows.}
\label{tab:headline}
\begin{tabular}{lrrr}
\toprule
Model & Full & Liu ($n{=}9{,}175$) & Kim ($n{=}11{,}334$) \\
\midrule
\textbf{PE-RankFormer (final)} & \textbf{0.9079} & \textbf{0.8585} & \textbf{0.8124} \\
PE-RankFormer (previous round) & 0.8933 & 0.8462 & 0.7836 \\
PE-RankFormer (first version) & 0.8865 & 0.8349 & 0.7751 \\
OptiPrime & 0.8690 & 0.8365 & 0.7320 \\
\bottomrule
\end{tabular}
\end{table}

\begin{table}[H]
\centering
\caption{Paired protospacer-clustered bootstrap, 5{,}000 resamples, on identical rows.
Clustering is by protospacer because pegRNAs sharing a target are not independent; a
row-level bootstrap would materially understate the intervals.}
\label{tab:boot}
\begin{tabular}{lrlrr}
\toprule
Comparison & $\Delta\rho$ & 95\% CI & wins & $p$ \\
\midrule
\textbf{Final $-$ OptiPrime} & $\mathbf{+0.0389}$ & $[+0.0286, +0.0498]$ & 1.000 & $<0.0002$ \\
Final $-$ previous round & $+0.0146$ & $[+0.0113, +0.0184]$ & 1.000 & $<0.0002$ \\
Previous round $-$ OptiPrime & $+0.0243$ & $[+0.0138, +0.0346]$ & 1.000 & $<0.0002$ \\
\bottomrule
\end{tabular}
\end{table}

\paragraph{A caveat that belongs beside the headline.}
The absolute figure's own 95\% interval is $[0.8955, 0.9173]$, and $90.0\%$ of
resamples exceed $0.90$. The $0.90$ threshold is therefore cleared by the point
estimate but \emph{not} established at 95\% confidence. The \emph{margin over
OptiPrime} is a much stronger claim than the absolute level: it wins every one of
5{,}000 resamples.

% ==========================================================================
\section{Data and evaluation protocol}
% ==========================================================================

\subsection{Corpus}

All models are trained on OptiPrime's own training mix, supplied by its authors:
297{,}962 rows, plus the 20{,}509-row held-out set, giving 318{,}471 total. Three
source studies contribute, and their proportions differ sharply between training and
evaluation --- a fact that becomes important in \S\ref{sec:negatives}.

\begin{table}[H]
\centering
\caption{Corpus composition. Schwank contributes 58.4\% of training rows and 0\% of the
evaluation set. ``Zero-mass'' is the fraction of rows with efficiency exactly $0$.}
\label{tab:corpus}
\begin{tabular}{lrrr}
\toprule
Source & \% of training & \% of held-out & Zero-mass \\
\midrule
Schwank (\texttt{pridict\_pridict2}) & 58.4\% & 0.0\% & 10.3\% \\
Liu (\texttt{hsu2026}) & 22.0\% & 44.7\% & 0.8\% \\
Kim (\texttt{deepprime}) & 19.6\% & 55.3\% & 49.9\% \\
\bottomrule
\end{tabular}
\end{table}

\subsection{Evaluation discipline}

A four-level hierarchy separates exploration from confirmation:

\begin{enumerate}
\item \textbf{Matched development folds} (3 folds, $\approx$35.7k rows each), composed
to match the held-out set's Liu/Kim proportions. Used freely for exploration.
\item \textbf{Internal lockbox} (17{,}975 rows, 367 protospacers, sharing no protospacer
with any development validation set). Touched once, to gate a shortlist.
\item \textbf{Official 5-fold out-of-fold predictions.} Every row is scored only by the
checkpoint that held its fold out.
\item \textbf{Official held-out set.} Evaluated once, after the model was frozen in
writing, with the expected result recorded in advance.
\end{enumerate}

Everything reported as a development number is out-of-fold. The selection protocol
(which subset, which combination rule, how to break ties) was written down and
committed \emph{before} the ensemble search ran.

% ==========================================================================
\section{Model architecture}
% ==========================================================================

\subsection{Backbone}

Each member is a dual-encoder relational model:

\begin{enumerate}
\item an \textbf{edit encoder} over paired wild-type/edited tokens;
\item a \textbf{pegRNA encoder} over segment-aware spacer/PBS/RTT tokens;
\item \textbf{bidirectional cross-attention} between the two streams;
\item \textbf{FiLM conditioning} on experimental context (cell line, PE system, Cas9
variant), applied either once after pooling (``late'') or at every block (``layerwise'');
\item an \textbf{attention-pooled output head}.
\end{enumerate}

Base configuration: $d_{\text{model}}{=}384$, 6 edit layers, 4 pegRNA layers, 2 cross
blocks, $\approx$25M parameters.

\subsection{The ordinal outcome head}

The evaluation metric is Spearman, which depends only on order. The ordinal head
targets that directly. Fix $K-1$ thresholds $t_k$ at quantiles of the \emph{training}
efficiency distribution; the head emits one logit per threshold, trained with binary
cross-entropy against the cumulative indicators, and scores a row by

\[
s(x) \;=\; \frac{1}{K-1}\sum_{k=1}^{K-1}\sigma(z_k(x))
\;\approx\; \text{the row's normalised rank.}
\]

Two properties motivate it. It is \emph{metric-matched}: a head estimating the target's
quantile optimises the evaluated quantity, whereas a proportion head is only
incidentally monotone. And it is \emph{robust to the target's skew} --- efficiency has
large mass at zero, so a Huber or proportion objective spends most of its gradient on a
few high-efficiency rows, while $K-1$ balanced binary decisions spread supervision
across the distribution.

Thresholds are computed from training rows only and stored in the checkpoint, so
evaluation reproduces the score mapping exactly.

\subsection{The state-space sequence mixer}

One family of members replaces self-attention with a \textbf{bidirectional diagonal
state-space (S4D) mixer}. With a diagonal state matrix the convolution kernel is
available in closed form,
\[
K[\ell] \;=\; 2\,\mathrm{Re}\!\left(\textstyle\sum_n C_n \frac{e^{\Delta A_n}-1}{A_n}\,
\left(e^{\Delta A_n}\right)^{\ell}\right),
\]
and is applied by FFT convolution. Two independent kernels (forward and time-reversed)
are summed, since prime-editing efficiency depends on context on both sides of the
edit. The block topology deliberately mirrors the Transformer layer it replaces, so a
comparison isolates the mixing mechanism rather than confounding it with block design.

\subsection{Which axis matters: an objective\,$\times$\,architecture factorial}

All four cells existed as out-of-fold predictions, so this was computed from stored
artifacts at no additional cost.

\begin{table}[H]
\centering
\caption{Out-of-fold development Spearman, fold-averaged.}
\label{tab:factorial}
\begin{tabular}{lrr}
\toprule
 & Simplex head & Ordinal head \\
\midrule
Transformer & 0.8798 & 0.8911 \\
State-space (SSM) & 0.9011 & \textbf{0.9082} \\
\bottomrule
\end{tabular}
\end{table}

\begin{center}
\begin{tabular}{lr}
\toprule
Main effect of \textbf{objective} (ordinal $-$ simplex) & $+0.0092$ \\
Main effect of \textbf{architecture} (SSM $-$ Transformer) & $\mathbf{+0.0192}$ \\
Interaction & $-0.0043$ \\
\bottomrule
\end{tabular}
\end{center}

\paragraph{This corrects an earlier claim of ours.} An earlier version of this work
described the ordinal head as ``the breakthrough $\ldots$ a change of objective, not
architecture''. The factorial does not support that: the architecture effect is roughly
twice the objective effect. The ordinal head was the more novel contribution and the
one that produced decorrelated errors, but the state-space mixer contributed more raw
accuracy.

The interaction is \emph{sub-additive}: ordinal supervision is worth $+0.0114$ on a
Transformer but only $+0.0071$ on an SSM. The two changes partly capture the same
improvement rather than stacking, which sets a realistic ceiling on
``add another mechanism'' strategies --- a prediction borne out in \S\ref{sec:negatives}.

Caveat: these are observational cells rather than a purpose-built factorial. Three are
from one training round; the Transformer/simplex cell is an earlier frozen baseline.
They share corpus, splits, out-of-fold scoring, capacity and optimiser recipe, so the
main effects are trustworthy at the $\pm0.002$ development resolution, but the
interaction term is the quantity most exposed to that mismatch.

% ==========================================================================
\section{The final system}
% ==========================================================================

\subsection{Composition}

An equal-weight \textbf{rank average} of five out-of-fold members. Each member is five
checkpoints, one per official fold; every row is scored only by the checkpoint that
held its fold out. Rank averaging is used because members are not mutually calibrated
--- the ordinal members emit rank estimates and the simplex members emit proportions,
so averaging raw scores would mostly measure scale mismatch.

\begin{table}[H]
\centering
\caption{Members and their individual held-out performance.}
\label{tab:members}
\begin{tabular}{llrrr}
\toprule
Member & Design & Full & Liu & Kim \\
\midrule
\texttt{ordSSM} & ordinal head $+$ S4D mixer & \textbf{0.9082} & 0.8576 & \textbf{0.8151} \\
\texttt{ssm} & simplex head $+$ S4D mixer & 0.9017 & 0.8563 & 0.7982 \\
\texttt{ordA} & ordinal head $+$ layerwise FiLM & 0.9005 & 0.8394 & 0.8082 \\
\texttt{ordC} & ordinal head $+$ feature branch & 0.8947 & 0.8261 & 0.8027 \\
\texttt{familyA} & simplex head $+$ layerwise FiLM & 0.8888 & 0.8412 & 0.7787 \\
\midrule
\textbf{Ensemble} & equal-weight rank average & \textbf{0.9079} & \textbf{0.8585} & 0.8124 \\
\bottomrule
\end{tabular}
\end{table}

\subsection{An honest observation about the ensemble}
\label{sec:ens-honest}

Table~\ref{tab:members} contains a result we did not anticipate and think should be
stated plainly: \textbf{the single best member matches the five-member ensemble on the
held-out set} ($0.9082$ vs $0.9079$). A leave-one-out analysis on held-out data makes
the point sharper.

\begin{table}[H]
\centering
\caption{Leave-one-member-out on the held-out set. Positive contribution means the
ensemble is worse without that member. \emph{Post-hoc analysis: this did not and could
not inform the frozen model choice.}}
\label{tab:loo}
\begin{tabular}{lrr}
\toprule
Member dropped & Ensemble without it & Contribution \\
\midrule
\texttt{ordSSM} & 0.9048 & $+0.0031$ \\
\texttt{ssm} & 0.9066 & $+0.0013$ \\
\texttt{ordA} & 0.9073 & $+0.0006$ \\
\texttt{ordC} & 0.9085 & $-0.0007$ \\
\texttt{familyA} & 0.9092 & $-0.0013$ \\
\bottomrule
\end{tabular}
\end{table}

On the development folds the ensemble was worth $+0.0074$ over its best member
($0.9156$ vs $0.9082$). On held-out that advantage is $-0.0003$. \textbf{The ensemble
gain did not transfer.} Two members would, post-hoc, have been marginally better
excluded.

We report this rather than quietly presenting the ensemble as the better system,
because it bears on how the result should be used. All the differences involved
($0.0003$--$0.0013$) are far below the $\approx$0.005 resolution of our development
evaluation, so the honest reading is that \textbf{the ensemble and its best single
member are indistinguishable on this benchmark}. For practical use, the single
ordinal-SSM model is preferable: one model instead of twenty-five checkpoints, at the
same accuracy. We did not re-select on held-out to make that switch official, because a
model chosen by inspecting the test set carries none of the guarantees the frozen one
does.

\subsection{Selection procedure}

The ensemble was chosen under a protocol committed in writing before the search ran:
rank subsets on development folds; require a candidate to beat the incumbent on
\emph{every} fold rather than on the mean; gate the shortlist on the lockbox, scored
out-of-fold; and on disagreement prefer the lockbox winner, breaking near-ties toward
fewer members.

The development winner (4 members) and the lockbox winner (5 members) disagreed by
$0.001011$, eleven millionths past the pre-set $0.001$ near-tie threshold, so the rule
selected the five-member subset. We note this was marginal --- the two candidates are
statistically indistinguishable --- and followed the rule as written rather than
re-arguing it after seeing the numbers. The lockbox's own unconstrained best subset
($0.9155$) was \emph{not} adopted, since selecting on the lockbox would have converted
the only unworn evaluation surface into a fitted one.

% ==========================================================================
\section{Detailed results}
% ==========================================================================

\subsection{Per-condition breakdown}

\begin{table}[H]
\centering
\caption{Held-out Spearman by experimental condition ($n \geq 300$). Performance spans
$0.21$ between the weakest and strongest conditions.}
\label{tab:percond}
\begin{tabular}{lllrr}
\toprule
Study & Cell line & PE system & $n$ & $\rho$ \\
\midrule
Kim & A549 & PE2 & 1299 & 0.6712 \\
Kim & MDA-MB-231 & PE2 & 620 & 0.7100 \\
Kim & A549 & PE4 & 1281 & 0.7144 \\
Kim & HCT116 & PE2 & 651 & 0.7866 \\
Kim & NIH3T3 & PE4 & 650 & 0.8034 \\
Kim & HeLa & PE2 & 664 & 0.8103 \\
Liu & HEK293T & PE2 & 2158 & 0.8179 \\
Liu & HEK293T & PE4 & 2148 & 0.8282 \\
Kim & HEK293T & PE2 & 3208 & 0.8298 \\
Liu & HeLa & PE4 & 2435 & 0.8311 \\
Liu & HeLa & PE2 & 2434 & 0.8318 \\
Kim & HEK293T & PE4 & 1135 & 0.8592 \\
Kim & DLD1 & PE4 & 1164 & 0.8778 \\
Kim & DLD1 & PE2 & 662 & 0.8809 \\
\bottomrule
\end{tabular}
\end{table}

Performance tracks the zero-mass of the condition: A549/PE2 is 63.9\% exact zeros and
scores $0.671$; DLD1/PE2 is 37.8\% zeros and scores $0.881$.

\subsection{Calibration: restoring an absolute prediction}

Rank averaging produces excellent orderings but destroys the output scale: the
ensemble's ``prediction'' is a normalised rank, not an editing efficiency, so absolute
error metrics against it are meaningless. We fit a monotone map
$g:\text{rank}\mapsto\text{efficiency}$ using isotonic regression on out-of-fold
development predictions only, and apply it once to the frozen held-out predictions.

\begin{table}[H]
\centering
\caption{Effect of monotone calibration. The map is fitted entirely on development
data, is verified non-decreasing, and inverts no pair --- it is a change of units on an
already-frozen prediction vector, not a new model selection.}
\label{tab:calib}
\begin{tabular}{lrrrr}
\toprule
 & Spearman & Pearson & MAE & RMSE \\
\midrule
Uncalibrated & 0.9079 & 0.7278 & 0.3869 & 0.4343 \\
\textbf{Calibrated} & \textbf{0.9079} & \textbf{0.8637} & \textbf{0.0478} & \textbf{0.0912} \\
\bottomrule
\end{tabular}
\end{table}

MAE falls roughly eight-fold and Pearson rises $0.136$ with the ranking untouched (the
residual $2\times10^{-5}$ Spearman shift comes from isotonic regression's
piecewise-constant plateaus creating ties, not from any pair being reordered). The
system now emits a usable efficiency estimate.

% ==========================================================================
\section{What did not work}
\label{sec:negatives}
% ==========================================================================

After the system above was frozen, approximately fifty further experiments were run
across three rounds. \textbf{None improved on it.} We report them because the pattern
is more informative than any individual result, and because several are ideas a reader
might otherwise expect to see tried.

\begin{table}[H]
\centering
\caption{Summary of failed interventions. ``$\Delta$'' is against a matched control ---
in every case an unmodified re-run of the same model, which we found to be
indispensable (see below).}
\label{tab:negatives}
\begin{tabular}{llr}
\toprule
Category & Intervention & $\Delta$ \\
\midrule
Post-hoc combination
 & context-gated ensemble weights & $+0.0000$ \\
 & nonlinear stacking (GBM, ridge) & $+0.0013$ \\
 & residual learning on features & $+0.0006$ (oracle bound) \\
\midrule
Supervision geometry
 & auxiliary simplex head (dual-head) & $+0.0008$ \\
 & multi-resolution ordinal & $+0.0008$ \\
 & context-relative ordinal (auxiliary) & $+0.0016$ \\
 & quantile head (pinball loss) & $-0.0140$ \\
 & pairwise ranking loss & $-0.0037$ \\
\midrule
Output parameterisation
 & true CORAL (rank-consistent) & $+0.0005$ \\
 & monotonicity penalty & $+0.0007$ \\
 & hurdle head (zero-inflation) & $-0.0022$ \\
\midrule
Architecture / capacity
 & hybrid SSM$+$attention (alternating) & $-0.0014$ \\
 & hybrid SSM$+$attention (parallel gate) & $-0.0027$ \\
 & SSM state dim 128 / 256 & $-0.0022$ / $+0.0002$ \\
 & wider FFN, deeper stack & $-0.0001$ / did not replicate \\
 & mixture-of-experts (4 / 8) & did not replicate / $-0.0013$ \\
 & layerwise context on SSM & $+0.0001$ \\
\midrule
Training scheme
 & protospacer bagging & $+0.0002$ to $-0.0022$ \\
 & source re-weighting (mild / aggressive) & did not replicate / $-0.0059$ \\
 & new random seed & $+0.0022$ \\
\bottomrule
\end{tabular}
\end{table}

\subsection{Three findings from the failures}

\paragraph{A mechanism-free control is indispensable.}
In one round, fourteen candidates each showed a small positive gain of $+0.0017$ to
$+0.0020$. The round also included a plain re-run of the baseline at a different batch
size --- no new mechanism whatsoever --- which gained $+0.0018$. \textbf{The difference
attributable to every idea in that round combined was $+0.0002$.} Without that control
we would have reported fourteen promising leads.

\paragraph{Single-fold results at this effect size are not evidence.}
In a later round, three interventions on genuinely orthogonal axes (capacity, data
distribution, conditional capacity) each reached $+0.003$ to $+0.004$ on one
development fold. All three \emph{changed sign} on a second fold. Replication, not
effect size, was the discriminating test.

\paragraph{Some benchmark gains are metric artifacts.}
28.4\% of rows have efficiency exactly zero (49.9\% within Kim). Collapsing the
lowest-scoring predictions onto a single tied value matches that tie structure and
gains $+0.0044$ pooled Spearman, positive on all three folds. We \emph{rejected} it:
applying the same transform to five other models gave $+0.0127$ to $+0.0162$
regardless of model quality, so it exploits Spearman's tie convention rather than
improving prediction, and OptiPrime emits continuous tie-free scores. This is worth
recording as a benchmark caveat --- \textbf{comparisons between models with different
tie behaviour on zero-inflated data are not automatically like-for-like.}

% ==========================================================================
\section{Remaining headroom, and where it sits}
% ==========================================================================

The obvious inference from \S\ref{sec:negatives} --- that the benchmark is saturated
--- is testable, and it is wrong.

\subsection{A replicate-based noise ceiling}

The Kim data contains genuine technical replicates: 649 groups in which all sixteen
design and condition covariates \emph{and} the target site are identical. If the same
pegRNA in the same condition is measured twice and the measurements disagree, no model
can predict either better than they predict each other. For two independent
measurements of a true value, the reliability $R = \mathrm{corr}(y_1,y_2)$ implies a
ceiling of $\sqrt{R}$ against a single observation.

Two corrections were necessary, and either shortcut inverts the conclusion:

\begin{itemize}
\item \textbf{Naive replicate keys give $R=0.7987$}, almost exactly our model's Kim
score --- which would have ``confirmed'' saturation. But grouping on design and
condition alone merges epegRNA with plain-pegRNA constructs, which differ in motif,
epegRNA flag and linker. Those are different designs, not repeats.
\item \textbf{Treating an observed zero as noiseless inflates the ceiling.} 21.4\% of
rows measured at exactly $0$ have a non-zero replicate (median $0.0015$). A zero is
\emph{censored}, not certain. A Gaussian noise model assuming otherwise gives $0.968$;
an empirical nearest-neighbour resampler gives $0.9026$, and the difference falls
almost entirely on the zero-heavy conditions the model does worst on.
\end{itemize}

\begin{table}[H]
\centering
\caption{Kim headroom against the empirical ceiling, development fold.}
\label{tab:ceiling}
\begin{tabular}{lrrr}
\toprule
Scope & Model & Ceiling & Gap \\
\midrule
All Kim rows & 0.7869 & 0.9026 & $\mathbf{+0.1157}$ \\
Rows that edit ($y>0$) & 0.8465 & 0.9489 & $+0.1024$ \\
\bottomrule
\end{tabular}
\end{table}

\textbf{Kim is not noise-limited}, and the gap is not merely the zero block: excluding
zeros entirely still leaves $+0.10$.

\subsection{Diagnosis: context rescales, it does not reorder}

For designs measured in two conditions, across 45 condition pairs:

\begin{center}
\begin{tabular}{lr}
\toprule
Mean cross-condition $\rho$, \textbf{true} (what the biology does) & 0.683 \\
Mean cross-condition $\rho$, \textbf{model} (what it predicts) & 0.835 \\
Excess & $\mathbf{+0.152}$ \\
\bottomrule
\end{tabular}
\end{center}

The model applies a near-universal sequence ranking and uses experimental context
mainly to set the scale, whereas reality reorders designs substantially between cell
lines and PE systems. \textbf{That reordering is learnable}: predicting the
cross-condition rank shift from just sixteen engineered design features reaches
held-out $\rho \approx 0.275$. The excess similarity ($+0.152$) is the same order as
the gap to the ceiling ($+0.116$).

\subsection{A falsified fix, and what it taught}

The natural remedy --- conditioning at every block instead of once after pooling ---
was implemented for the state-space backbone and tested on three folds with matched
controls. It was null ($+0.0001$), and the pre-registered mechanism test failed in the
informative direction: cross-condition rank correlation \emph{rose} to $0.875$ from the
control's $0.828$.

More conditioning made the model \emph{more} condition-invariant. FiLM is structurally
a per-channel scale and shift, so applying it at every block adds no capacity to
reorder designs; it adds a better pathway for the cross-condition mean shift, which
explains far more pooled variance and is much easier to fit. The extra capacity is
spent on the shortcut.

The diagnosis survives; the instrument was wrong. Work is ongoing on removing the
shortcut from the \emph{objective} rather than adding conditioning capacity.

% ==========================================================================
\section{Limitations}
% ==========================================================================

\begin{enumerate}
\item \textbf{The $0.90$ threshold is not established at 95\% confidence}
($[0.8955, 0.9173]$; 90\% of resamples above). The margin over OptiPrime is the robust
claim.
\item \textbf{The ensemble does not beat its best member on held-out}
(\S\ref{sec:ens-honest}). The development-fold advantage did not transfer.
\item \textbf{196 of 649 Kim replicate groups straddle the official train/held-out
split}, so those held-out rows have an exact design-and-condition twin in training.
This affects any model trained on this split equally --- ours and OptiPrime's --- so it
does not distort the comparison, but it makes absolute held-out scores slightly
optimistic.
\item \textbf{Kim has no fresh evaluation surface left.} Of 2{,}501 Liu+Kim training
protospacers, 1{,}553 sit in development validation sets, 367 were consumed by the
lockbox, and 302 collide with Schwank; the 279 remaining are 98.5\% Liu. Kim-side
claims rest on out-of-fold development folds and a once-used lockbox.
\item \textbf{Our OptiPrime reproduction scores above its published figure} under every
aggregation convention we tried, and no evaluation script was released. The direction
favours the baseline, so it does not undermine our claims, but it caps the precision of
the comparison.
\item \textbf{No external validation.} All results are on one benchmark.
\end{enumerate}

% ==========================================================================
\section{Reproduction}
% ==========================================================================

\begin{verbatim}
# Final system: five members x five official folds
scripts/train/train_pilot.py --config configs/round5/baseline_ordinal_ssm.yaml \
    --run-name <member> --val-fold <k> [member-specific flags]

# Out-of-fold development predictions
scripts/evaluate/run_round4_phase3.sh

# Ensemble search and lockbox gate
scripts/evaluate/search_ensemble.py --members ...

# Held-out evaluation (guarded; every run is audit-logged)
scripts/evaluate/evaluate_heterogeneous_heldout.py --member ... \
    --allow-heldout-evaluation
scripts/evaluate/round4_final_bootstrap.py

# Calibration and noise ceiling
scripts/evaluate/calibrate_ensemble.py --apply-heldout
scripts/evaluate/noise_ceiling_empirical.py
\end{verbatim}

The frozen model specification, the pre-registered selection protocol, and the full
research log including every negative result are in \texttt{reports/}.

\end{document}
